\begin{frame}{왜 속도를 2배·100배·수천 배로 늘리기가 이렇게 어려운가}
  \begin{itemize}
    \item 현대 CPU: 8$\sim$64개 코어 --- 코어가 16개면
      \textbf{독립 작업} 16개까지 동시 가능
    \item 그런데 ``환경 하나를 스텝시키기''는 \textbf{독립 작업이
      아니다}
  \end{itemize}
  \vskip0.5em
  \begin{block}{\kb{시간 방향 의존성}}
    한 환경의 $t$스텝은 같은 환경의 $t-1$스텝 결과가 나와야만
    계산이 시작된다.
  \end{block}
  \vskip0.5em
  \begin{itemize}
    \item 이 의존성을 깨는 방법은 ``같은 환경의 시간을 잘게
      나누는 것''이 \textbf{아니다}
    \item $\Rightarrow$ \textbf{서로 다른 환경}(=초깃값이 다른
      복제본)을 코어마다 하나씩 맡기는 것
    \item 예: 16코어 CPU로 Ant 16개를 병렬 $\Rightarrow$
      Ant 16개 환경 스텝을 동시 계산 --- 이것이 ``CPU 병렬''의
      \textbf{전부}
  \end{itemize}
\end{frame}
